\section{Curry--Howard correspondence}
\label{sec:normalization-curry-howard-03}

Under the Curry--Howard correspondence, a proposition is represented by a
type, a derivation by a typed term, implication introduction by lambda
abstraction, implication elimination by application, and normalization by
computation.  The correspondence is structural: it does not identify every
formal derivation with a program without first fixing a calculus and a
term assignment.

\begin{example}[Composition]
The term
\[
  \lambda f{:}(B\to C).\,
  \lambda g{:}(A\to B).\,
  \lambda x{:}A.\,f(gx)
\]
has type
\[
  (B\to C)\to(A\to B)\to A\to C.
\]
Reading types as propositions gives a proof of the corresponding intuitionistic
formula.  Each abstraction discharges one typed assumption, and each
application is an implication-elimination step.
\end{example}

Products and sums correspond to conjunction and disjunction after their term
constructors, eliminators, and reduction rules are defined.  The empty type
corresponds to \(\bottom\).  Classical principles require additional control
operators, translations, or axioms; they are not terms of the pure simply
typed lambda calculus.

